% 国赛论文骨架（全国大学生数学建模竞赛论文格式规范，2026 年修订稿）
% 编译：XeLaTeX（中文）。电子版构建时注释掉“承诺书/编号专用页”两节，
% 第一页必须为摘要专用页；纸质版取消注释并打印装订。
\documentclass[a4paper,12pt]{ctexart}

% ---- 第一条：A4，页边距上下左右 ≥2.5cm，左侧装订 ----
\usepackage[top=2.5cm,bottom=2.5cm,left=2.5cm,right=2.5cm]{geometry}
\usepackage{fancyhdr}
\usepackage{graphicx}
\usepackage{booktabs}
\usepackage{amsmath}
\usepackage{hyperref}

% ---- 第三条：页码从摘要页起，页脚中部，阿拉伯数字从 1 连续编号 ----
\fancypagestyle{mainmatter}{%
  \fancyhf{}\fancyfoot[C]{\thepage}\renewcommand{\headrulewidth}{0pt}}
\newcommand{\startmainnumbering}{\clearpage\setcounter{page}{1}\pagestyle{mainmatter}}

% ---- 第六条/第十一条：全文匿名。提交前 grep 学校、姓名、赛区 ----
% ANONYMITY-CHECK: grep -r "大学\|学院\|赛区" 正文 附录 支撑材料

\begin{document}

% ===== 纸质版专用（电子版构建时注释掉下面两节） =====
% ---- 第二条：第一页承诺书 ----
\section*{承诺书}
% 按当年官方专用页原样填入。

% ---- 第二条：第二页编号专用页 ----
\section*{编号专用页}
% 按当年官方专用页原样填入。

% ===== 电子版第一页必须为摘要专用页 =====
% ---- 第三条：摘要 ≤1 页，含标题和关键词，无需英文 ----
\startmainnumbering
\section*{摘要}
% 关键词：XXX；XXX；XXX

% ---- 第四条：正文不要目录，不超过 30 页 ----
\section{问题重述}
\section{问题分析}
\section{模型假设}
\section{符号说明}
\section{模型建立与求解}
\section{结果分析与检验}
\section{模型的评价与推广}
\section{参考文献}
% ---- 第七条：引用他人或公开资料（含网上资料）文内标注 + 文末规范列出 ----
\begin{thebibliography}{99}
\bibitem{ref1} % 作者. 题名[文献类型]. 出版地: 出版者, 年.
\end{thebibliography}

% ---- 第四条/第五条：附录页数不限，打印装订一起交 ----
\appendix
\section{附录：支撑材料文件列表}
% 将支撑材料 RAR/ZIP 包内文件逐个列出（第十一条）。
% - 程序源文件（含 EXCEL/SPSS 交互命令），须完整可运行，结果与论文相符
% - 自主查阅的数据资料（赛题原始数据除外）
% - 大篇幅中间结果图表
\begin{itemize}
  \item % 文件名 —— 说明
\end{itemize}

\section{附录：源程序代码}
% 二选一，删其一：
% （有程序）下面粘贴全部完整、可运行源程序：
% \begin{verbatim}
% \end{verbatim}
% （无程序）本论文没有用到程序。
% （无支撑材料时另注）本论文没有支撑材料。

\end{document}
